\chapter{Proposta Metodológica}
\label{ch:proposta}

% Meta editorial: 22--32 páginas. Este é o capítulo central. A ordem narrativa
% é problema -> formulação -> decisões de projeto -> diagnósticos -> validação.

\section{Visão geral da ANT}
% Responder antes das equações: qual problema resolve; que similaridades usa;
% quais relações deixam de produzir pressão; quais continuam ativas; onde a
% ANT entra na destilação contrastiva do TagFex. Inserir diagrama comparando
% TagFex e TagFex+ANT.

\section{Atualizações contrastivas não essenciais}

\subsection{Cenário motivador e interpretação geométrica}
% Exemplo visual pequeno com âncora, positivo, negativos adequados e violadores.
% Mostrar intuitivamente por que relações adequadas não precisam receber a
% mesma pressão de otimização indefinidamente.

\subsection{Definição operacional e consequências esperadas}
% Distinguir loss residual, gradiente não nulo, violação de margem, atualização
% não essencial e dano ao conhecimento anterior. Definir quais serão observáveis
% e quais permanecem inferências a testar.

\subsection{Predições e diagnósticos do mecanismo}
% Se a hipótese estiver correta: adjusted loss, active ratio, distribuição de
% violações e hard-negative similarity devem evoluir de forma compatível, sem
% perda de plasticidade medida no desempenho das novas classes.

\section{Formulação base da ANT}
% Desenvolver em sequência, com interpretação após cada equação:
% (1) batch e duas views; (2) embeddings normalizados; (3) matriz de
% similaridade; (4) máscaras de self/positivos; (5) negativos válidos;
% (6) referência; (7) margem; (8) v_ij = s_ij - r + gamma;
% (9) termos ativos; (10) agregação; (11) L_ANT;
% (12) alpha L_InfoNCE + beta L_ANT; (13) fluxo de gradientes.
% Conferir tudo contra _build_ant_matrix(), _compute_ant_loss() e
% _compute_contrastive_loss_base() antes da redação definitiva.

\section{Decisão de projeto I: referência da ANT}

\subsection{Referência global (aGlobal)}
% Um valor compartilhado extraído do domínio válido da matriz ANT. Formalizar
% domínio, acoplamento entre âncoras, margem efetiva e gradientes esperados.

\subsection{Referência por âncora (aLocal)}
% Máximo por linha. Formalizar a diferença para aGlobal e discutir adaptação a
% graus distintos de dificuldade, sem declarar benefício antes das ablações.

% Inserir tabela: variante | domínio da referência | granularidade | efeito
% esperado | hipótese/experimento.

\section{Decisão de projeto II: cobertura da matriz}

\subsection{Cobertura intra-view assimétrica}
% Variante original: bloco Q1/S11, B âncoras e B-1 negativos. Explicitar quais
% views e blocos não participam.

\subsection{Cobertura simétrica completa (aSymFull)}
% Matriz 2B x 2B, removendo self e par positivo; 2B-2 negativos por âncora.
% Explicar simetria, cobertura, custo e comportamento esperado.

% Figura obrigatória: Q1, Q2, Q3, Q4, self, positivos, negativos e blocos usados
% por cada variante. Tabela: referência | cobertura | granularidade | custo.

\section{Decisão associada: normalização da InfoNCE}

\subsection{Normalização nGlobal e nLocal}
% Apresentar as duas formas lado a lado, por linha, com matriz pequena. Tratar
% como decisão da InfoNCE, separada da referência ANT.

\subsection{Equivalência algébrica e verificação numérica}
% A implementação nLocal subtrai uma constante por linha antes de uma expressão
% invariável a deslocamentos. Demonstrar condições de equivalência, comparar
% valores/gradientes e só então discutir estabilidade de ponto flutuante.
% TODO crítico: resolver a inconsistência entre a narrativa histórica de
% gradientes balanceados e a implementação atual.

\subsection{Combinações das decisões}
% aGlobal+nGlobal, aGlobal+nLocal, aLocal+nGlobal, aLocal+nLocal e aSymFull com
% a normalização correspondente. Explicar que são produtos de decisões
% independentes, não cinco métodos autônomos.

\section{Formulações alternativas como ablações}
% Consolidar logsumexp, expm1, softplus, top-k e active-only em uma tabela:
% motivação | diferença matemática | comportamento esperado | hipótese |
% experimento/status. Não elevar alternativas sem evidência a contribuições.
% Explicar aqui o piso log(N_valid) da logsumexp e a diferença entre problema
% diagnóstico e problema de otimização.

\section{Investigações auxiliares}

\subsection{Avg-K: estabilização temporal do professor}
% Média dos últimos K checkpoints de ta_net/projector; relação complementar ou
% independente da ANT; K=3/5/10; custo e evidência disponível. Não apresentar
% como contribuição principal sem suporte adicional.

\subsection{SBS: seleção de exemplares por velocidade de aprendizagem}
% Score por acertos/tentativas, filtragem das caudas, herding posterior,
% relação com rehearsal e limitações atuais (FFCV/distribuído).

\section{Algoritmo integrado}

\subsection{TagFex baseline e TagFex+ANT}
% Pseudocódigo comparável: forward -> professor -> matriz -> InfoNCE ->
% referência/margem/máscara ANT -> combinação -> backward -> atualização ->
% memória -> alinhamento -> avaliação. Destacar visualmente as linhas novas.

\subsection{Integração opcional dos mecanismos auxiliares}
% Mostrar Avg-K/SBS como módulos opcionais após o núcleo, não como parte
% necessária da definição da ANT. Incluir complexidade de tempo/memória.

\section{Metodologia experimental}

\subsection{Matriz de hipóteses e experimentos}
% Tabela obrigatória: pergunta/hipótese | comparação controlada | datasets |
% seeds/status | métricas de desempenho | métricas de mecanismo.

\subsection{Protocolos e reprodutibilidade}
% Datasets e divisões; classes iniciais/incrementais; seeds; arquitetura;
% épocas; batch e views; memória; hiperparâmetros; baselines; hardware;
% configuração e origem do resultado. Marcar completo/parcial/em execução/falha.

\subsection{Métricas e análise estatística}
% Avg Acc@1, NME@1, desempenho final, forgetting e curvas; raw/adjusted loss,
% active ratio, hard-negative similarity, violações/percentis e gaps; média,
% dispersão, tamanho de efeito e testes compatíveis com o desenho final.

\section{Resultados preliminares e análise do mecanismo}

\subsection{Comportamento interno da ANT}
% Estrutura de cada resultado: hipótese -> expectativa -> métrica -> figura ->
% observação -> interpretação -> limitação. Explicar matematicamente o piso:
% log(254) ~= 5,54; raw loss quase plana não implica ausência de aprendizado.
% Analisar adjusted loss, active ratio, hard negatives e percentis.

\subsection{Desempenho incremental e ablações}
% Accuracy/NME/forgetting com múltiplas seeds quando disponíveis; separar
% resultados completos, parciais, em execução, falhos e exploratórios.
% Organizar por decisão: referência, cobertura, normalização, formulação;
% Avg-K/SBS vêm depois como investigações auxiliares.

\subsection{Síntese das hipóteses}
% Classificar evidência como forte, tendência, inconclusiva ou não sustentada.
% Não converter diferenças pequenas de uma seed em superioridade.

\section{Ameaças à validade e limitações}
% Em uma seção coesa: validade interna, externa e de construto; seeds;
% hiperparâmetros; baseline; datasets/backbones; custo; dependência do TagFex;
% relação entre proxies geométricas e atualização realmente não essencial.

\section{Extensões em investigação}
% Separar da contribuição defendida: colisões entre tarefas, negativos baseados
% em protótipos antigos e outras estratégias somente após validação diagnóstica.
